% =====================================================================
% 研究ノート（正本）。使い方はリポジトリ直下の AGENTS.md。
% プロジェクト: meta_contrib（メタ研究。DECISIONS.md D15）
% ロック前は、確証用の標本で格と特徴量を突き合わせない（data/README.md）。
% =====================================================================
\documentclass[a4paper,10pt]{ltjsarticle}
% 共通プリアンブル。<project>/paper/main.tex から % 共通プリアンブル。<project>/paper/main.tex から % 共通プリアンブル。<project>/paper/main.tex から \input{../../common/lab_preamble}
% 箱は決まったものだけ使う（AGENTS.md §3）。増やすときは AGENTS.md の一覧も直す。
\usepackage{amsmath,amssymb,booktabs,tabularx,longtable,enumitem,xurl,graphicx}
\usepackage[margin=25mm]{geometry}
\usepackage{xcolor}
\usepackage[most]{tcolorbox}
\usepackage{tikz,pgfplots}
\pgfplotsset{compat=1.18}
\usetikzlibrary{arrows.meta,decorations.pathreplacing,positioning,calc}
\usepackage[backend=biber,style=authoryear,natbib=true,maxcitenames=2,doi=true,url=true]{biblatex}
\usepackage[hidelinks]{hyperref}

\definecolor{labBlue}{HTML}{1F5FA8}
\definecolor{labGreen}{HTML}{2E7D4F}
\definecolor{labPurple}{HTML}{6A4C9C}
\definecolor{labOrange}{HTML}{C8641E}
\definecolor{labRed}{HTML}{B03A3A}
\definecolor{labGray}{HTML}{5A6270}

\tcbset{labbase/.style={breakable,enhanced,boxrule=0.6pt,arc=2pt,left=6pt,right=6pt,top=4pt,bottom=4pt,
  fonttitle=\bfseries\small,before skip=8pt,after skip=8pt}}

% 読むための箱
\newtcolorbox{chaptermap}{labbase,colback=labBlue!4,colframe=labBlue,title={この文書の地図}}
\newtcolorbox{keypoint}{labbase,colback=labBlue!4,colframe=labBlue,title={要点}}
\newtcolorbox{term}[1]{labbase,colback=labGray!5,colframe=labGray,title={用語：#1}}
\newtcolorbox{pitfall}{labbase,colback=labRed!4,colframe=labRed,title={落とし穴}}
\newtcolorbox{debate}{labbase,colback=labGray!5,colframe=labGray,title={論争・別の説明}}

% 研究の箱
% 研究仮説：\begin{hypothesis}{H1 短い名前}
\newtcolorbox{hypothesis}[1]{labbase,colback=labGreen!5,colframe=labGreen,title={仮説：#1}}
% 識別戦略：\begin{design}{H1 差の差}
\newtcolorbox{design}[1]{labbase,colback=labBlue!4,colframe=labBlue,title={識別戦略：#1}}
% 推定結果：\begin{result}{H1 の検証（2026-11、確証）}
\newtcolorbox{result}[1]{labbase,colback=labPurple!4,colframe=labPurple,title={結果：#1}}
% 意思決定：続ける／やめる／変える
\newtcolorbox{decision}{labbase,colback=labOrange!6,colframe=labOrange,title={判断（Go/No-Go）}}

% 箱は決まったものだけ使う（AGENTS.md §3）。増やすときは AGENTS.md の一覧も直す。
\usepackage{amsmath,amssymb,booktabs,tabularx,longtable,enumitem,xurl,graphicx}
\usepackage[margin=25mm]{geometry}
\usepackage{xcolor}
\usepackage[most]{tcolorbox}
\usepackage{tikz,pgfplots}
\pgfplotsset{compat=1.18}
\usetikzlibrary{arrows.meta,decorations.pathreplacing,positioning,calc}
\usepackage[backend=biber,style=authoryear,natbib=true,maxcitenames=2,doi=true,url=true]{biblatex}
\usepackage[hidelinks]{hyperref}

\definecolor{labBlue}{HTML}{1F5FA8}
\definecolor{labGreen}{HTML}{2E7D4F}
\definecolor{labPurple}{HTML}{6A4C9C}
\definecolor{labOrange}{HTML}{C8641E}
\definecolor{labRed}{HTML}{B03A3A}
\definecolor{labGray}{HTML}{5A6270}

\tcbset{labbase/.style={breakable,enhanced,boxrule=0.6pt,arc=2pt,left=6pt,right=6pt,top=4pt,bottom=4pt,
  fonttitle=\bfseries\small,before skip=8pt,after skip=8pt}}

% 読むための箱
\newtcolorbox{chaptermap}{labbase,colback=labBlue!4,colframe=labBlue,title={この文書の地図}}
\newtcolorbox{keypoint}{labbase,colback=labBlue!4,colframe=labBlue,title={要点}}
\newtcolorbox{term}[1]{labbase,colback=labGray!5,colframe=labGray,title={用語：#1}}
\newtcolorbox{pitfall}{labbase,colback=labRed!4,colframe=labRed,title={落とし穴}}
\newtcolorbox{debate}{labbase,colback=labGray!5,colframe=labGray,title={論争・別の説明}}

% 研究の箱
% 研究仮説：\begin{hypothesis}{H1 短い名前}
\newtcolorbox{hypothesis}[1]{labbase,colback=labGreen!5,colframe=labGreen,title={仮説：#1}}
% 識別戦略：\begin{design}{H1 差の差}
\newtcolorbox{design}[1]{labbase,colback=labBlue!4,colframe=labBlue,title={識別戦略：#1}}
% 推定結果：\begin{result}{H1 の検証（2026-11、確証）}
\newtcolorbox{result}[1]{labbase,colback=labPurple!4,colframe=labPurple,title={結果：#1}}
% 意思決定：続ける／やめる／変える
\newtcolorbox{decision}{labbase,colback=labOrange!6,colframe=labOrange,title={判断（Go/No-Go）}}

% 箱は決まったものだけ使う（AGENTS.md §3）。増やすときは AGENTS.md の一覧も直す。
\usepackage{amsmath,amssymb,booktabs,tabularx,longtable,enumitem,xurl,graphicx}
\usepackage[margin=25mm]{geometry}
\usepackage{xcolor}
\usepackage[most]{tcolorbox}
\usepackage{tikz,pgfplots}
\pgfplotsset{compat=1.18}
\usetikzlibrary{arrows.meta,decorations.pathreplacing,positioning,calc}
\usepackage[backend=biber,style=authoryear,natbib=true,maxcitenames=2,doi=true,url=true]{biblatex}
\usepackage[hidelinks]{hyperref}

\definecolor{labBlue}{HTML}{1F5FA8}
\definecolor{labGreen}{HTML}{2E7D4F}
\definecolor{labPurple}{HTML}{6A4C9C}
\definecolor{labOrange}{HTML}{C8641E}
\definecolor{labRed}{HTML}{B03A3A}
\definecolor{labGray}{HTML}{5A6270}

\tcbset{labbase/.style={breakable,enhanced,boxrule=0.6pt,arc=2pt,left=6pt,right=6pt,top=4pt,bottom=4pt,
  fonttitle=\bfseries\small,before skip=8pt,after skip=8pt}}

% 読むための箱
\newtcolorbox{chaptermap}{labbase,colback=labBlue!4,colframe=labBlue,title={この文書の地図}}
\newtcolorbox{keypoint}{labbase,colback=labBlue!4,colframe=labBlue,title={要点}}
\newtcolorbox{term}[1]{labbase,colback=labGray!5,colframe=labGray,title={用語：#1}}
\newtcolorbox{pitfall}{labbase,colback=labRed!4,colframe=labRed,title={落とし穴}}
\newtcolorbox{debate}{labbase,colback=labGray!5,colframe=labGray,title={論争・別の説明}}

% 研究の箱
% 研究仮説：\begin{hypothesis}{H1 短い名前}
\newtcolorbox{hypothesis}[1]{labbase,colback=labGreen!5,colframe=labGreen,title={仮説：#1}}
% 識別戦略：\begin{design}{H1 差の差}
\newtcolorbox{design}[1]{labbase,colback=labBlue!4,colframe=labBlue,title={識別戦略：#1}}
% 推定結果：\begin{result}{H1 の検証（2026-11、確証）}
\newtcolorbox{result}[1]{labbase,colback=labPurple!4,colframe=labPurple,title={結果：#1}}
% 意思決定：続ける／やめる／変える
\newtcolorbox{decision}{labbase,colback=labOrange!6,colframe=labOrange,title={判断（Go/No-Go）}}

\addbibresource{../common/refs.bib}

\title{一流誌を分けるのは派手さではなく、やり方を変える力である（仮説）\\計量経済学の雑誌の格と「渋い問い」のメタ研究}
\author{（担当者名）}
\date{2026年11月24日（最終更新）}

\begin{document}
\maketitle

\begin{abstract}
経済学では、見た目の渋い問いが一流誌に載り、広く引かれることが多い。
本研究は、計量経済学の一流・二流・三流誌の論文の題目と要旨を、貢献の型と書き方のはい／いいえの質問で判定し、どの型が格を分けるかを調べる。
主な仮説は、一流誌を分けるのは題材の派手さではなく、多くの研究者のやり方や信念を変える貢献（誤りを正す・道具を作る・まとめる枠組み）だ、というものである。
仮説の候補27個から、人が10個を検証に選んだ。すべて未検証・未ロックである。
次に、雑誌の格付け表と判定の質問を作り、データを取る。
\end{abstract}

\begin{chaptermap}
\textbf{3行でいうと}
\begin{enumerate}[leftmargin=*]
  \item 背景：IF を目的変数にすると、分野の被引用の癖（サーベイや一部の出版社が高く、理論の雑誌が低い）を拾ってしまう。雑誌の格で測り直す。
  \item 方法：題目と要旨に、貢献の型と書き方の質問を決定モデルで当てる。格ごとの該当割合を比べ、内容と文体を分ける。
  \item 現状：候補27個から10個を検証に選んだ（H01, H02, H03, H07, H11, H16, H18, H22, H23, H25）。未検証・未ロック。標本は探索用30\%・確証用70\%に分ける予定。
\end{enumerate}
\end{chaptermap}

%======================================================================
\section{背景：「渋い問い」はなぜ評価されるのか}\label{sec:bg}
経済学の一流誌には、一般の人には細かい技術の話に見える論文が多く載る。
例えば、差の差法の標準誤差が過小になりうることを示した論文\parencite{bertrand2004}、
処置の時期がずれるときの二方向固定効果の推定値の解釈を示した論文\parencite{goodmanbacon2021,dechaisemartin2020,callaway2021}、
弱い操作変数の問題を扱った論文\parencite{staiger1997}、
効果がばらつくときに操作変数推定が何を推定するかを示した論文\parencite{imbens1994}である。
これらに共通するのは、表面は小さくても、読んだ後にやり方や信念を変える研究者が多いことだ、というのが筆者の見解（仮説）である。
ただし、以上の論文はいずれも書誌だけを照合しており、本文は確かめていない。

題目と要旨から掲載誌の IF を予測した予備的な分析（非公開）が、この研究の着想のもとになった。
一方で IF は、サーベイや被引用の多い分野の雑誌で高く、理論の雑誌で低い。
このため IF を目的変数にすると、論文の質ではなく分野の被引用の慣行を学んでしまうおそれがある。
本研究は、目的変数を外部の格付けリストによる雑誌の格に替えて、貢献の型と書き方を分けて調べる。

\begin{term}{雑誌の格}
一流・二流・三流の3段階。通念ではなく、外部の格付けリスト（ABS/AJG、ABDC）で決める。リストで格が割れる雑誌は感度分析に回す。
\end{term}

\begin{term}{決定モデル}
文章ではなく、確率・選択・点数を返す LLM。本研究では、はい／いいえの質問に対する P(はい) を特徴量に使う。
\end{term}

%======================================================================
\section{先行研究}\label{sec:prior}
未調査。経済学の雑誌の格付けと、論文の特徴と掲載先・被引用の関係を調べた研究を、文献マスター表に足してから書く。
結果を見る前に決める意味と限界は \textcite{olken2015}、有意な結果に偏って報告される問題は \textcite{brodeur2020} を参照。

%======================================================================
\section{研究仮説：10個を検証する}\label{sec:hyp}
候補27個は \texttt{hypotheses/hypotheses.json} にある。2026-11-24 に、人が次の10個を検証に選んだ（◆③。H22 は同日に追加）：H01, H02, H03, H07, H11, H16, H18, H22, H23, H25。
成功基準と棄却条件の数値は仮置きで、ロック前に人が決める。
H22（文体を消しても判定が変わらないか）と H23（判定と人の一致）は、他の仮説すべての前提となる測定の点検である。H25（「一流誌らしい文体」の分解）は格を使わないので、ロック前に全標本で実行してよい。

\begin{table}[htbp]
\centering\small
\begin{tabularx}{\linewidth}{llXl}
\toprule
群 & ID & 仮説（短い名前） & 検証 \\
\midrule
貢献の型 & H01 & 誤りを正す型は一流誌に多い & ◎ \\
 & H02 & 小幅な拡張は三流誌に多い & ◎ \\
 & H03 & 道具を作る型は一流誌に多い & ◎ \\
 & H04 & まとめる枠組みの型は一流誌に多い &  \\
 & H05 & 帰無の結果の書き方が格で違う &  \\
 & H06 & 一流誌の要旨は貢献を一つに絞る &  \\
一般性と射程 & H07 & 特定の国・業種の名前は三流誌に多い & ◎ \\
 & H08 & 一流誌の手法論文は広い使い手に向けて書く &  \\
 & H09 & 一流誌の手法論文には実証例が付く &  \\
 & H10 & 一流誌の手法はモデルの族に対して主張する &  \\
証拠の種類 & H11 & シミュレーションだけの手法論文は三流誌に多い & ◎ \\
 & H12 & 理論の保証は格を上げる（IF とは逆） &  \\
書き方 & H13 & 「有意」で語る応用論文は三流誌に多い &  \\
 & H14 & 一流誌の応用論文は効果の大きさを数値で書く &  \\
 & H15 & 一流誌の要旨は「問題→解決」で始まる &  \\
 & H16 & 流行語は三流誌に多い（渋さの裏返し） & ◎ \\
 & H17 & 決まり文句の要旨は三流誌に多い &  \\
内容と文体 & H18 & 貢献の型は文体を超えて格を説明する & ◎ \\
 & H19 & 同じ雑誌の中でも、貢献の型は被引用を予測する &  \\
時間と拡散 & H20 & 話題は一流誌から三流誌へ遅れて広がる &  \\
 & H21 & 誤りを正す型は一流誌で増えている &  \\
測定の妥当性 & H22 & 判定は文体を消しても変わらない & ◎ \\
 & H23 & 決定モデルの判定は人の判定と一致する & ◎ \\
 & H24 & 格を目的変数にすると、理論の向きが IF と逆になる &  \\
「一流誌らしい文体」の分解 & H25 & 「一流誌らしい文体」は具体的な文体の特徴で説明できる & ◎ \\
 & H26 & 書き換え実験：一流誌の文体に直すと何が変わるか &  \\
 & H27 & 分解した文体の特徴でも、雑誌の格を当てられる &  \\
\bottomrule
\end{tabularx}
\caption{仮説の候補27個と、◆③で検証に選んだ10個（◎）。選外の仮説も消さずに残す。}
\label{tab:mc-hyp}
\end{table}

\begin{hypothesis}{H01 誤りを正す型は一流誌に多い}
\textbf{仮説}：「誤りを正す」型の割合は、一流誌で三流誌の2倍以上である。\\
\textbf{推定対象}：一流誌と三流誌の、該当割合の比。\\
\textbf{成功基準}：比が2以上で、比の95\%信頼区間（雑誌単位のクラスタ・ブートストラップ）が1を含まない。\\
\textbf{棄却条件}：比の95\%信頼区間の上限が1.2を下回る。\\
\textbf{否定されたとき}：誤りを正す型は格を分けない。「渋い問い」の説明は別に要る。\\
\textbf{状態}：未検証（未ロック）
\end{hypothesis}

\begin{hypothesis}{H02 小幅な拡張は三流誌に多い}
\textbf{仮説}：「小幅な拡張」型の割合は、三流誌で一流誌の2倍以上である。\\
\textbf{推定対象}：三流誌と一流誌の、該当割合の比。\\
\textbf{成功基準}：比が2以上で、比の95\%信頼区間（雑誌単位のクラスタ・ブートストラップ）が1を含まない。\\
\textbf{棄却条件}：比の95\%信頼区間の上限が1.2を下回る。\\
\textbf{否定されたとき}：小幅な拡張は格を分けない。\\
\textbf{状態}：未検証（未ロック）
\end{hypothesis}

\begin{hypothesis}{H03 道具を作る型は一流誌に多い}
\textbf{仮説}：「汎用の道具」型の割合は、一流誌で三流誌の1.5倍以上である。\\
\textbf{推定対象}：一流誌と三流誌の、該当割合の比。\\
\textbf{成功基準}：比が1.5以上で、比の95\%信頼区間（雑誌単位のクラスタ・ブートストラップ）が1を含まない。\\
\textbf{棄却条件}：比の95\%信頼区間の上限が1.1を下回る。\\
\textbf{否定されたとき}：道具を作ることは格を分けない。\\
\textbf{状態}：未検証（未ロック）
\end{hypothesis}

\begin{hypothesis}{H07 特定の国・業種の名前は三流誌に多い}
\textbf{仮説}：応用論文のうち、要旨に特定の国・地域・業種が出る割合は、三流誌で一流誌の2倍以上である。\\
\textbf{推定対象}：応用論文に限った、該当割合の比。\\
\textbf{成功基準}：比が2以上で、比の95\%信頼区間（雑誌単位のクラスタ・ブートストラップ）が1を含まない。\\
\textbf{棄却条件}：比の95\%信頼区間の上限が1.2を下回る。\\
\textbf{否定されたとき}：対象の狭さは格を分けない。一般性は書き方で補えている。\\
\textbf{状態}：未検証（未ロック）
\end{hypothesis}

\begin{hypothesis}{H11 シミュレーションだけの手法論文は三流誌に多い}
\textbf{仮説}：手法論文のうち、証拠がシミュレーションだけのものの割合は、三流誌で一流誌の2倍以上である。\\
\textbf{推定対象}：手法論文に限った、該当割合の比。\\
\textbf{成功基準}：比が2以上で、比の95\%信頼区間（雑誌単位のクラスタ・ブートストラップ）が1を含まない。\\
\textbf{棄却条件}：比の95\%信頼区間の上限が1.2を下回る。\\
\textbf{否定されたとき}：証拠の種類は格を分けない。\\
\textbf{状態}：未検証（未ロック）
\end{hypothesis}

\begin{hypothesis}{H16 流行語は三流誌に多い（渋さの裏返し）}
\textbf{仮説}：流行の題材を題目か要旨の1文目に出す割合は、三流誌で一流誌の1.5倍以上である。\\
\textbf{推定対象}：該当割合の比。\\
\textbf{成功基準}：比が1.5以上で、比の95\%信頼区間（雑誌単位のクラスタ・ブートストラップ）が1を含まない。\\
\textbf{棄却条件}：比の95\%信頼区間の上限が1.1を下回る。\\
\textbf{否定されたとき}：流行語は格を分けない。「渋さ」は題材の選び方ではない。\\
\textbf{状態}：未検証（未ロック）
\end{hypothesis}

\begin{hypothesis}{H18 貢献の型は文体を超えて格を説明する}
\textbf{仮説}：文体の質問だけのモデルに貢献の型の質問を足すと、格の予測の順位相関（論文単位、交差検証）が0.05以上上がる。\\
\textbf{推定対象}：交差検証での順位相関の上がり幅。\\
\textbf{成功基準}：上がり幅が0.05以上で、95\%信頼区間が0を含まない。\\
\textbf{棄却条件}：上がり幅の95\%信頼区間の上限が0.02を下回る。\\
\textbf{否定されたとき}：格の差は文体でほぼ説明できる。内容の差は文体に現れている。\\
\textbf{状態}：未検証（未ロック）
\end{hypothesis}

\begin{hypothesis}{H22 判定は文体を消しても変わらない}
\textbf{仮説}：要旨を LLM で中立な文体に言い換えても、内容の質問の答えは元の答えと相関0.8以上で一致する。言い換え後の要旨で、格の予測の順位相関の下がり幅は0.05以下である。\\
\textbf{推定対象}：内容の質問ごとの、言い換え前後の P(はい) の相関。格の予測の順位相関の変化。\\
\textbf{成功基準}：内容の質問の相関の中央値が0.8以上で、予測の下がり幅が0.05以下。\\
\textbf{棄却条件}：内容の質問の相関の中央値が0.6を下回るか、予測の下がり幅が0.15を超える。\\
\textbf{否定されたとき}：内容の判定に文体の手がかりが混ざっている。H01〜H18 は文体の効果を含む可能性があると書く。\\
\textbf{状態}：未検証（未ロック）
\end{hypothesis}

\begin{hypothesis}{H23 決定モデルの判定は人の判定と一致する}
\textbf{仮説}：人が付けた80本（無作為60本＋語で選んだ20本）の正解に対して、型ごと（論文の種類・誤りを正す・小幅な拡張・道具・国や業種・流行の題材・証拠の種類）の AUC の中央値が0.8以上である。\\
\textbf{推定対象}：型ごとの AUC・正解率・コーエンの κ（無作為60本と80本全体）。\\
\textbf{成功基準}：型ごとの AUC の中央値が0.8以上（無作為60本で）。\\
\textbf{棄却条件}：型ごとの AUC の中央値が0.7を下回る。\\
\textbf{手順}：\texttt{research/labeling\_protocol.md}。人が付けるのは型の単位で、1本あたり6〜9個の判断。約2時間半。\\
\textbf{状態}：未検証（未ロック）
\end{hypothesis}

\begin{hypothesis}{H25 「一流誌らしい文体」は具体的な文体の特徴で説明できる}
\textbf{仮説}：具体的な文体の質問20〜30問（書き出しの型、貢献の明示、主張の範囲の限定、流行語、決まり文句、文法、長さ、能動態など）で、「一流誌のような文体ですか」の P(はい) のばらつきの70\%以上を説明できる。\\
\textbf{推定対象}：具体的な文体の質問から「一流誌らしさ」の P(はい) を予測するモデルの、交差検証での決定係数。質問ごとの寄与（Shapley 値）。\\
\textbf{成功基準}：交差検証の決定係数が0.7以上。\\
\textbf{棄却条件}：交差検証の決定係数が0.4を下回る（判定は具体的な書き方以外のもの、例えば雑誌や論文そのものの知識に頼っている）。\\
\textbf{否定されたとき}：「一流誌らしい文体」は書き方の束では説明できない。LLM の事前知識の可能性が高まる（H26 で確かめる）。\\
\textbf{状態}：未検証（未ロック）
\end{hypothesis}

%======================================================================
\section{データ}\label{sec:data}
予定。詳細は \texttt{data/README.md}。
格は ABDC Journal Quality List 2025 で定義する（人の判断、2026-11-24）。一流は A*（Econometrica、Journal of Econometrics、Econometric Theory、JBES、Quantitative Economics、Review of Economics and Statistics、Journal of Applied Econometrics、Econometrics Journal、Journal of Financial Econometrics の9誌）、二流は A（Econometric Reviews、Journal of Time Series Analysis、Oxford Bulletin of Economics and Statistics、International Journal of Forecasting、Empirical Economics、Applied Economics、Journal of Forecasting、Studies in Nonlinear Dynamics \& Econometrics の8誌）、三流は B（Econometrics〈MDPI〉、Econometrics and Statistics、Journal of Econometric Methods、Journal of Time Series Econometrics、Journal of Quantitative Economics、Applied Economics Letters の6誌）である。
当初の仮の3段階（通念による）は、ABDC と大きく食い違った。仮の3段階は感度分析に回す。一覧は \texttt{research/journal\_tiers.csv}。
OpenAlex から、2015〜2024年の論文を各誌最大300本取る。
秘匿区分は T0 である。要旨の本文は再配布しない。
OpenAlex では一部の雑誌で要旨が大きく欠ける（Journal of Econometrics など、2015〜2024年の論文の15〜19\%しか要旨が無い）。Semantic Scholar で補い、残る欠けを雑誌×年で記録する（\texttt{data/README.md}）。

判定の質問の案は \texttt{questions/questions.toml} にある（50問。論文の種類3問、仮説ごとに3〜4問、「一流誌らしい文体」1問、具体的な文体24問）。人が確かめてからロックで固定する。

%======================================================================
\section{推定の方法}\label{sec:design}
\begin{design}{格ごとの該当割合の比}
\textbf{比べるもの}：質問ごとに、格ごとの該当割合とその比（一流対三流）。主な分析は格を目的変数にした順序ロジット。\\
\textbf{仮定}：(a) 決定モデルの判定が人の判定と一致する（H23 で確かめる）。(b) 判定が雑誌の手がかりや文体に引きずられない（H22 で確かめる。雑誌名を伏せる・1誌ずつ抜く点検も置く）。(c) 雑誌の格は外部のリストで決まる。\\
\textbf{反証テスト}：要旨を中立な文体に言い換えて判定し直す。雑誌名を伏せる。Econometric Theory など1誌ずつ抜いて結果が変わらないかを見る。\\
\textbf{推論}：雑誌単位のクラスタ・ブートストラップ。格ごとの雑誌数が6〜8と少ないので、ワイルド・ブートストラップも併記する。多重検定の扱い（主要仮説の指定か補正）はロック前に決める。\\
\textbf{探索と確証}：標本を雑誌ごとに層を切って探索用30\%・確証用70\%に分ける。ロック前に格と特徴量を突き合わせてよいのは探索用だけである。
\end{design}

%======================================================================
\section{予想される結果}\label{sec:expected}
未作成。雑誌の数が少ないので、検出力は雑誌単位のばらつきで決まる。標本を決めてから \texttt{sim/} で出す。

%======================================================================
\section{結果}\label{sec:result}
\begin{result}{全仮説（未実施）}
\textbf{状態}：未検証。データは未取得。
\end{result}

%======================================================================
\section{考察}\label{sec:discussion}
\begin{pitfall}
要旨は著者が書いた自己申告である。一流誌の要旨に「誤りを正す」と書く癖があるだけなら、H01 は内容ではなく書き方の差を測っている。
H22（要旨を中立な文体に言い換えて判定し直す）で、この別の説明を確かめる。掲載後の被引用で確かめる H19 は選外にしたので、その分は残る懸念として扱う。
\end{pitfall}

%======================================================================
\section{判断と次の一手}\label{sec:decision}
\begin{decision}
\textbf{結論}：プロジェクトとして始める（人の判断、2026-11-24。DECISIONS.md D15）。\\
\textbf{決めたこと（2026-11-24）}：判定には OpenAI Decisions API を使う（利用の許可あり）。雑誌の並びは\ref{sec:data}節の案でよい（格付けリストでの照合は行う）。\\
\textbf{決めたこと（2026-11-24、◆③）}：H01, H02, H03, H07, H11, H16, H18, H22, H23, H25 を検証する（H22 は同日に追加）。選んだ理由は未記入。\\
\textbf{決めたこと（2026-11-24）}：雑誌の格は ABDC 2025 で定義し、三流（B）の雑誌を足して23誌にする。\\
\textbf{次の一手}：(1) 人が質問案（\texttt{questions/questions.toml}）と正解づけの手順書（\texttt{research/labeling\_protocol.md}）を確かめる。(2) AI が要旨の欠けを補い（\texttt{code/01b\_fill\_gaps.py}）、標本を探索用と確証用に分ける。(3) 人が正解づけ用の80本に答えを付ける（H23）。(4) AI が探索用の標本で判定を回す。
\end{decision}

%======================================================================
\section*{更新履歴}
\begin{itemize}
  \item 2026-11-24：作成。仮説の候補24個を出した。
  \item 2026-11-24：判定モデル（OpenAI Decisions API）と雑誌の並びを人が決めた。「一流誌らしい文体」を分解する仮説 H25〜H27 を足した。
  \item 2026-11-24：◆③ 検証する仮説を9個選んだ（人の判断）。同日、H22 を追加して10個にした。
  \item 2026-11-24：正解づけの手順書を作り、H23 の正解を「100本・質問ごと」から「80本・型ごと」に改めた（ロック前）。
  \item 2026-11-24：雑誌の格を ABDC 2025 で定義し直した（人の判断）。当初の仮の3段階は ABDC と大きく食い違ったため。三流の雑誌を5誌足して23誌にした。
\end{itemize}

\printbibliography[title={参考文献}]
\end{document}
